
\chapter{مسیر یادگیری Python و Django}

\section{هدف این فصل}

چون Backend قبلی توسعه‌دهنده \eng{Node.js/Express} بوده، این فصل ترتیب یادگیری را طوری می‌چیند که دانش قبلی حفظ شود و تفاوت‌های Python/Django به‌صورت مرحله‌ای فهمیده شوند.

\section{Python: چیزهایی که واقعاً لازم داریم}

\subsection{مرحله P1: Syntax و Data Types}

باید بتوانی این موارد را بدون کمک توضیح بدهی و بنویسی:
\begin{itemize}
  \item variable؛
  \item \eng{str/int/float/bool/None}؛
  \item list؛
  \item tuple؛
  \item dict؛
  \item set؛
  \item if/elif/else؛
  \item for/while.
\end{itemize}

مثال:
\begin{latin}
\begin{lstlisting}[language=Python]
topics = ["science", "history", "business"]

for topic in topics:
    print(topic)
\end{lstlisting}
\end{latin}

\subsection{مرحله P2: Function}

\begin{latin}
\begin{lstlisting}[language=Python]
def calculate_reading_minutes(word_count: int) -> int:
    words_per_minute = 220
    return max(1, round(word_count / words_per_minute))
\end{lstlisting}
\end{latin}

مفاهیم:
\begin{itemize}
  \item parameter؛
  \item argument؛
  \item return؛
  \item default parameter؛
  \item type hint؛
  \item scope.
\end{itemize}

\subsection{مرحله P3: Module و Package}

در Node.js با import/export آشنا هستی. در Python نیز Code را به Moduleها تقسیم می‌کنیم.

\begin{latin}
\begin{lstlisting}[language=Python]
from recommendations.scoring import calculate_score
\end{lstlisting}
\end{latin}

باید تفاوت این‌ها را بفهمی:
\begin{itemize}
  \item file/module؛
  \item package؛
  \item import path؛
  \item \filepath{__init__.py}.
\end{itemize}

\subsection{مرحله P4: Class و Object}

Django Modelها Class هستند؛ بنابراین قبل از Model باید OOP پایه را بفهمی.

\begin{latin}
\begin{lstlisting}[language=Python]
class IdeaSummary:
    def __init__(self, title: str, score: float):
        self.title = title
        self.score = score

    def is_recommended(self) -> bool:
        return self.score >= 5
\end{lstlisting}
\end{latin}

مفاهیم:
\begin{itemize}
  \item class؛
  \item instance؛
  \item attribute؛
  \item method؛
  \item constructor؛
  \item inheritance؛
  \item composition.
\end{itemize}

\subsection{مرحله P5: Exception}

\begin{latin}
\begin{lstlisting}[language=Python]
try:
    score = calculate_score(user, idea)
except ValueError as exc:
    logger.warning("Invalid recommendation input: %s", exc)
\end{lstlisting}
\end{latin}

نباید هر Exception را با \tech{except Exception} پنهان کنیم.

\subsection{مرحله P6: Comprehension و Iterator}

در حد عملی:
\begin{latin}
\begin{lstlisting}[language=Python]
topic_ids = [topic.id for topic in idea.topics.all()]
\end{lstlisting}
\end{latin}

هدف فهم خوانایی و Lazy behavior است، نه استفاده نمایشی از Syntax پیچیده.

\section{Django از چه اجزایی تشکیل شده؟}

\subsection{Project}

\eng{Django Project} تنظیمات کل Application است.

فایل‌های اصلی:
\begin{itemize}
  \item \filepath{settings.py}
  \item \filepath{urls.py}
  \item \filepath{asgi.py}
  \item \filepath{wsgi.py}
\end{itemize}

\subsection{App}

Django App یک Domain/Capability قابل تفکیک است.

مثلاً:
\begin{latin}
\begin{lstlisting}
accounts
content
engagement
recommendations
\end{lstlisting}
\end{latin}

\begin{warningbox}
برای هر Table یک App نساز. App باید Business Responsibility داشته باشد.
\end{warningbox}

\section{URL Routing}

Express:
\begin{latin}
\begin{lstlisting}
app.get("/ideas/:id", handler)
\end{lstlisting}
\end{latin}

Django:
\begin{latin}
\begin{lstlisting}[language=Python]
path("ideas/<int:pk>/", views.idea_detail)
\end{lstlisting}
\end{latin}

DRF در مراحل بعد abstractionهای بیشتری می‌دهد، اما ابتدا باید این مسیر ساده فهمیده شود.

\section{Request و Response}

باید بفهمی:
\begin{itemize}
  \item Browser request می‌فرستد؛
  \item Django URL را Resolve می‌کند؛
  \item View اجرا می‌شود؛
  \item DB query در صورت نیاز انجام می‌شود؛
  \item Response برمی‌گردد.
\end{itemize}

\section{Model و ORM}

مثال آموزشی:
\begin{latin}
\begin{lstlisting}[language=Python]
from django.db import models

class Topic(models.Model):
    slug = models.SlugField(unique=True)
    name_fa = models.CharField(max_length=100)
    name_en = models.CharField(max_length=100)
\end{lstlisting}
\end{latin}

این Class به Table در Database نگاشت می‌شود.

باید این عملیات ORM را بلد باشی:
\begin{latin}
\begin{lstlisting}[language=Python]
Topic.objects.all()
Topic.objects.get(slug="science")
Topic.objects.filter(is_active=True)
Topic.objects.create(...)
\end{lstlisting}
\end{latin}

\section{Migration}

Migration تاریخچه تغییر Schema است.

Flow:
\begin{latin}
\begin{lstlisting}
change models.py
    ↓
python manage.py makemigrations
    ↓
inspect migration
    ↓
python manage.py migrate
\end{lstlisting}
\end{latin}

\begin{warningbox}
Migration را فایل Generated بی‌ارزش تصور نکن. بخشی از تاریخچه Schema پروژه است و باید Commit شود.
\end{warningbox}

\section{Django Admin}

Admin برای پروژه ما بسیار ارزشمند است چون قبل از ساخت Creator CMS می‌توانیم:
\begin{itemize}
  \item Topic بسازیم؛
  \item Creator ایجاد کنیم؛
  \item Source وارد کنیم؛
  \item Idea publish کنیم؛
  \item Reportها را بررسی کنیم.
\end{itemize}

بنابراین Data واقعی خیلی زود وارد Backend می‌شود.

\section{Custom User از روز اول}

\begin{warningbox}
تغییر \eng{AUTH\_USER\_MODEL} بعد از Migrationهای متعدد دردسر زیادی ایجاد می‌کند. بنابراین قبل از اولین Migration واقعی، Custom User را می‌سازیم.
\end{warningbox}

حتی اگر Custom User اولیه ساده باشد، انعطاف آینده را حفظ می‌کند.

\section{Django Authentication}

باید مفاهیم زیر را بفهمی:
\begin{itemize}
  \item password hashing؛
  \item authenticate؛
  \item login/logout؛
  \item session؛
  \item authenticated user؛
  \item group؛
  \item permission؛
  \item superuser.
\end{itemize}

\section{DRF بعد از Django پایه}

بعد از اینکه Model/View/URL را فهمیدی، می‌رویم سراغ:
\begin{itemize}
  \item Serializer؛
  \item validation؛
  \item APIView؛
  \item generic views؛
  \item ViewSet؛
  \item Router؛
  \item Permission؛
  \item Pagination.
\end{itemize}

\section{ترتیب تمرین‌های Backend}

\begin{enumerate}
  \item یک endpoint ساده health check؛
  \item Topic model + admin؛
  \item Source model؛
  \item Idea model؛
  \item query با shell؛
  \item Topic API read-only؛
  \item Idea API read-only؛
  \item create/update validation؛
  \item auth؛
  \item permissions؛
  \item engagement endpoints.
\end{enumerate}

\section{معیار فهم}

قبل از رفتن به مرحله بعد باید بتوانی پاسخ بدهی:
\begin{itemize}
  \item Model چه فرقی با Serializer دارد؟
  \item Migration دقیقاً چه چیزی را تغییر می‌دهد؟
  \item QuerySet چیست؟
  \item چرا \eng{ForeignKey} با IntegerField متفاوت است؟
  \item URLConf چه نقشی دارد؟
  \item Request چطور به View می‌رسد؟
  \item Permission کجا اعمال می‌شود؟
\end{itemize}
